\section{Introduction}
\label{sec:intro}

When one input to a multimodal system degrades, the usual evaluation asks
whether task performance falls. A drop has at least two causes. The
information needed for the task may have been destroyed. Or the information
may remain available in another input while the model keeps relying on the
degraded one. Accuracy alone does not distinguish these cases.

Substituting or removing an input to probe a model's dependence on it is
established practice (Section~\ref{sec:related}). What we study is a
severity-graded version: how performance and reliance change together as
one input is degraded step by step. For each severity we measure
\emph{utility}, the task performance with the degraded input relative to
removing it, and \emph{reliance}, the divergence between the model's
answer distributions with the degraded and with the removed input. We
locate where utility first reaches zero and ask how much reliance remains
there. To our knowledge, this pairing along a severity path has not been
studied in the work we reviewed.

\paragraph{Contributions.}
(i) We define the diagnostic: utility, reliance, a normalized reliance
$\hat R$, a utility-crossing severity, and a Misalignment Index (MI)
summarizing reliance beyond the crossing. (ii) We evaluate it on synthetic
fusion models with controlled reliance, including an oracle control. Our
first pre-set success criterion failed, and the revised criterion is
reported as post hoc. (iii) We apply it to three open VLMs on
multiple-choice VQA, with a held-out replication and an analysis of how
conclusions depend on the no-input baseline. (iv) We report an
over-reliance criterion pre-declared in a version-controlled log that is met in only
\checknum{2}{criterion:status:count_met} of 42 combinations. (v) We
include a small modality-dropout proxy study.
